\chapter{Logical View}
\ChapterLabel{LogicalView}
\section{Chapter Responsibility}
\label{Section:LogicalViewResponsibility}
Define the concepts, classes, interfaces, and responsibilities of the design.

\noindent\textbf{Expected result:} A structural model with explicit responsibility and ownership.

\section{How to Use This Chapter}
Model the domain concepts, responsibilities, relationships, data ownership, and invariants without committing to code-level structure. Use a class, entity, or conceptual diagram where useful. For every class and interface in a class diagram, list its primary responsibility in a titled and labeled register, as required by \hyperref[Requirement:DOCCLASS001]{DOC-CLASS-001}. Reference that register from the logical view and keep it synchronized with the diagram. Under \hyperref[Requirement:DOCLINK001]{DOC-LINK-001}, give each class or interface row a stable label and link class names in diagrams and prose to that row.
\section{Example}
\hyperref[Class:Request]{\textit{Request}} has an identifier, submitter, status, and submitted time. A request has one or more \hyperref[Class:RequestItem]{\textit{Request Items}}. Invariants: every identifier is unique; every request has a valid status; an approved request cannot return to Draft without an authorized transition.


Table~\ref{Table:ExampleClassResponsibilities} illustrates the required register.
\begin{tabularx}{\linewidth}{L{4.2cm}Y}
\caption{Example Class Responsibilities Register\label{Table:ExampleClassResponsibilities}}\\
\toprule Class & Primary responsibility\\\midrule
\phantomsection\label{Class:Request}\textit{Request} & Hold one request's identity, submitter, status, and submission time; enforce valid status transitions.\\
\phantomsection\label{Class:RequestItem}\textit{Request Item} & Hold one item belonging to a request and its item-specific data.\\\bottomrule
\end{tabularx}
